% Generated draft: requires mathematical and source review.
\documentclass[11pt]{article}
\usepackage[margin=25mm]{geometry}
\usepackage{fontspec}
\setmainfont{Latin Modern Roman}
\usepackage{amsmath,amssymb,mathtools,unicode-math}
\setmathfont{Latin Modern Math}
\usepackage{enumitem,needspace,xurl,hyperref}
\setlength{\parindent}{0pt}
\setlength{\parskip}{5pt}
\setlength{\emergencystretch}{3em}
\newcommand{\N}{\mathbb{N}}
\newcommand{\Z}{\mathbb{Z}}
\newcommand{\Q}{\mathbb{Q}}
\newcommand{\R}{\mathbb{R}}
\newcommand{\C}{\mathbb{C}}
\newcommand{\F}{\mathbb{F}}
\newcommand{\A}{\mathbb{A}}
\newcommand{\PP}{\mathbb{P}}
\newcommand{\E}{\mathbb{E}}
\newcommand{\eps}{\varepsilon}
\newcommand{\dd}{\,\mathrm{d}}
\begin{document}
\section*{UnsolvedMath ID: 11000063}
\textbf{Source problem code:} catalogue entry
\subsection{Definitions and mathematical statement}
The task is to give a new proof of the particular Hain theorem referred to in the preceding source text, using only elementary combinatorial-topological and group-theoretic methods, rather than the mixed-Hodge-theoretic proof described in the summary.

The supplied summary locates the theorem in the study of the Malcev completion of the Torelli group \(I_g\). In the standard closed-surface convention, \(I_g\) is the kernel of the mapping-class action on integral first homology; its rational Malcev, or unipotent, completion is the universal pro-unipotent group receiving its homomorphism over \(\Q\). However, the exact theorem, genus range, completion convention, generators, relations and hypotheses are not included in the extract. No particular presentation theorem is therefore asserted in its place.

The source hopes that such a proof would yield new understanding of \(I_g\) and help with a following problem to which mixed Hodge theory does not apply. That subsequent problem is missing; this motivation does not specify an additional theorem. Elementary here retains the explicit combinatorial-topological and group-theoretic restriction, rather than an invented difficulty standard.
\subsection{Short English statement}
Reprove the intended Hain theorem about Torelli-group Malcev completion by purely combinatorial, topological and group-theoretic methods; the theorem's precise statement is missing.
\subsection{Sources}
\begin{itemize}
\item B. Farb, chapter in \emph{Problems on Mapping Class Groups and Related Topics}, Problem 5.19, supplied PDF page 50. \url{https://www.math.uchicago.edu/~farb/papers/mcgbook.pdf}. Recorded access 2026-07-29; extraction pdf-font-structure.
\item Hain's theorem is named without a title, theorem number or full bibliography entry in this record. The supplied summary attributes its existing proof to Hodge theory and identifies the Malcev-completion setting.
\item Local UnsolvedMath registry, \texttt{data/unsolvedmath/problems.json}, permanent ID \(11000063\), exact source code \texttt{catalogue entry}; supplied OPEN-TRIAGE summary reports the elementary-proof challenge unresolved through 2026.
\end{itemize}
\end{document}
